\chapter{Wprowadzenie}

\textbf{Teoria automatów} stanowi jeden z fundamentalnych działów informatyki teoretycznej. Znajduje zastosowanie w matematycznym
modelowaniu systemów komputerowych, takich jak analiza leksykalna kompilatorów \cite[3.8.3 DFA's for Lexical Analyzers]{compilers},
czy \cite{}[TODO]. Zachowania takich systemów mogą zostać opisane w postaci formalnych modeli automatów, co umożliwia ich analizę oraz
wnioskowanie o własnościach przy użyciu narzędzi teorii języków formalnych.

\textbf{Aktywne uczenie się automatów} umożliwia stworzenie takiego modelu na podstawie obserwacji badanego systemu. Podejście to zakłada
dostęp do systemu w postaci czarnej skrzynki oraz wnioskowanie na podstawie jego odpowiedzi.

\textbf{Algorytm $L^*$ \cite{Angluin}}, opublikowany w roku 1987 przez Danę Angluin był jednym z pierwszych sukcesów w tej dziedzinie,
który przedstawił w jaki sposób można odtworzyć strukturę automatu na podstawie obserwacji. Zakładając możliwość odpowiadania na 2 rodzaje zapytań:
"Czy słowo $w$ jest akceptowane przez automat?" i "Czy dla zadanego automatu $H$, istnieje słowo które odróżnia go od uczonego systemu?",
można skonstruować minimalny deterministyczny automat skończony równoważny systemowi na podstawie którego odpowiadamy na zapytania.

\textbf{Automaty ze stosem} stanowią rozszerzenie automatów skończonych o dodatkową pamięć w postaci stosu. Umożliwiają reprezentowanie bardziej
skomplikowanych systemów, rozwiązujących szerszą gamę problemów, w szczególności te składające się z zagnieżdżonych struktur.
Jednakże ogólność tych automatów wiąże się z istotnymi ograniczeniami, jak nierozstrzygalność, lub koszt obliczeniowy wielu problemów decyzyjnych.

\textbf{Automaty Visibly Pushdown} stanowią klasę automatów zdefiniowaną w roku 2004 przez Rajeeva Alura i P.Madhusudana \cite{vpl}.
Jej obiekty podobnie jak automaty ze stosem, posiadają nieskończoną pamięć. Ich struktura jednakże jest bardziej przewidywalna,
co umożliwia wykonywanie m.in. takich operacji jak suma, różnica, czy przekrój dwóch automatów w czasie wielomianowym \cite{vpl}[2.3 Decision problems].

Automaty visibly pushdown znajdują zastosowanie w wielu systemach komputerowych. Przykładowo mogą służyć do przetwarzania dokumentów XML \cite{xml},
walidacji plików JSON \cite{json}, czy monitorowania programów w języku C \cite{Cmonitor}. Czyni je to wartościowymi obiektami badań.

\textbf{Aktywne uczenie się z dostępem do stosu.} W 2025 roku Jan Otop i Jakub Michaliszyn przedstawili algorytm do aktywnego uczenia się automatów DVPA.
Algorytm umożliwia poznanie dokładnej struktury automatu DVPA na podstawie obserwacji w wielomianowym czasie. Algorytm wymaga w porównaniu
do algorytmu $L^*$ dodatkowego zapytania o aktualny stan stosu, jak i możliwości modyfikacji stosu podczas wczytywania słowa przez automat.

Praca przedstawia implementację tego algorytmu oraz analizę jego wydajność pod kątem ilości wymaganych zapytań dla różnych klas automatów DVPA.
Badany jest również wpływ zastosowania wskazówki w postaci reguł przepisywania termów na podstawie pracy \textit{Active Automata Learning with Advice} \cite{srs},
której idea została dostosowana do automatów DVPA.

Pomimo szczegółowego opisu algorytmów w pracy \cite{main}, ich implementacja wymagała dogłębnego zrozumienia idei
aktywnego uczenia automatów, semantyki automatów visibly pushdown oraz opracowania rozbudowanych mechanizmów
weryfikacji poprawności. Wynikało to z konieczności zaprojektowania efektywnych struktur danych,
uzupełnienia brakujących szczegółów implementacyjnych oraz trudności w znajdywaniu powstałych błędów, wynikających z faktu,
że przewidywanie dokładnego zachowania algorytmu staje się skomplikowane nawet dla niewielkich automatów.